% =====================================================================
\chapter{C4: the tether-split key}\label{ch:c4}
% =====================================================================
\begin{keybox}[colframe=tether,colback=tether!5]
The TUAV's signing key is split in two, $\sk=a+g$. \textbf{The drone holds $a$; a small secure chip inside the ECRV holds $g$}, and the two are joined by
the tether. The drone can sign only when the chip adds its half over the wire. \textbf{If the tether is cut, the drone's half is useless}, and the ECRV
broadcasts a pre-signed \emph{kill notice}. Verifiers notice nothing: the public key and the signature are exactly as before.
\end{keybox}

\section{Why we add it: the TUAV is the new weak spot}
After C1 the TA has left the field, so the TUAV holds all the trust in its cell -- and it is the drone most worth capturing. The base paper never asks what
happens then.
\begin{center}
\small
\begin{tabularx}{\linewidth}{@{}l X@{}}
\toprule
\textbf{Design} & \textbf{What a thief obtains by capturing the TUAV}\\
\midrule
Base paper & $s_{tu}$, $\kappa_{tu}$ and every member's $\sigma_i$; beacons need only $s_{tu}$ (IS4) and key updates (DU1--DU2) are not signed, so the thief can
quietly keep control of the group key\\
C1 alone & the whole key $\sk_a$: it can run a fake TUAV from a laptop and a radio, signing beacons, handover replies and C3 votes\\
\textbf{C1 + C4} & \textbf{half a key, which signs nothing}, plus the current group key until the next rekey\\
\bottomrule
\end{tabularx}
\end{center}
\noindent While the drone is on the tether it hovers above the crew; it can fall into a stranger's hands only \emph{after} it leaves the tether -- a cut
cable, a power loss, an interceptor, or the ECRV driving away. C4 ties the key to exactly that moment.

\begin{figure}[H]
\centering
\begin{tikzpicture}[font=\sffamily\small]
  \pic[color=hoc,scale=1.7] at (0,3.3) {drone};
  \node[card=tether,anchor=west,text width=4.6cm] at (1.3,3.3) {\textbf{drone share} \;$a=x_A$\\ alone it signs nothing};
  \draw[tether,line width=1.6pt,decorate,decoration={snake,amplitude=0.6mm,segment length=4mm}] (0,2.65) -- (0,0.85);
  \node[font=\scriptsize\sffamily\bfseries,text=tether,align=left,anchor=west] at (0.25,1.75) {tether:\\ power + data};
  \pic[color=hoc,scale=1.8] at (0,0.35) {truck};
  \pic[color=vehc,scale=1.3] at (-1.6,0.45) {chip};
  \node[card=vehc,anchor=west,text width=4.6cm] at (1.3,0.35) {\textbf{chip share} \;$g=x_G+d$\\ in a SIM-sized secure chip\\ inside the ECRV};
  \node[card=deep,text width=5.0cm] at (10.3,1.85) {\textbf{public key unchanged}\\[2pt] $\pk=aP+gP=X+Y+h\Ppub$\\[2pt] nobody ever holds $a+g$};
  \draw[-{Stealth[length=3mm]},thick,deep] (6.15,3.3) -- (7.7,2.3);
  \draw[-{Stealth[length=3mm]},thick,deep] (6.15,0.35) -- (7.7,1.4);
\end{tikzpicture}
\caption{The split. Each share alone is useless; together they produce an ordinary signature under the TUAV's usual public key.}
\label{fig:split}
\end{figure}

\section{The construction}
\begin{stepbox}{tether}{S1--S4 \;Setup at the depot (once, before the mission)}
\steplabel{tether}{S1.} The drone picks $x_A$ and publishes $X_A=x_AP$; the chip picks $x_G$ and publishes $X_G=x_GP$ (each with a proof of possession);
$X=X_A+X_G$.\\
\steplabel{tether}{S2.} The TA registers the TUAV as in RG2 with $\ep=0$: $Y=yP$, \;$h=\Ho(\ID_a\cat0\cat X\cat Y)$, \;$d=y+s\,h$, and gives $d$ \textbf{to the
chip only}.\\
\steplabel{tether}{S3.} The chip keeps $g=x_G+d$; the drone keeps $a=x_A$.\\
\steplabel{tether}{S4.} Together they sign a \textbf{kill notice} $\mathit{KN}=$ ``$\ID_a$ is out of service'' and store it in the ECRV.
\end{stepbox}

\begin{stepbox}{tether}{J1--J5 \;Joint signing of a message $m$ over the tether (commit-then-reveal)}
\steplabel{tether}{J1.} Drone: pick $r_A$, $R_A=r_AP$; send $m$, $\ts$ and the \emph{commitment} $C_A=H(R_A)$ to the chip.\\
\steplabel{vehc}{J2.} Chip: refuse unless the tether is alive, $m$ is an allowed type (beacon, link offer, handover reply, C3 vote), $\ts$ is fresh and
increasing, and no other session is open. Pick $r_G$ and send $R_G=r_GP$.\\
\steplabel{tether}{J3.} Drone: send $R_A$. \;Chip: check $H(R_A)=C_A$.\\
\steplabel{vehc}{J4.} Both: $R=R_A+R_G$, \;$c=\Ht(m\cat\ID_a\cat0\cat X\cat Y\cat R\cat\ts)$. Chip sends $v_G=r_G+c\,g$ and logs $(m,\ts)$.\\
\steplabel{tether}{J5.} Drone: $v_A=r_A+c\,a$ and $v=v_A+v_G$. The signature $(R,v)$ has the usual form and size.
\end{stepbox}

\begin{figure}[H]
\centering
\resizebox{\linewidth}{!}{%
\begin{tikzpicture}[font=\sffamily\small]
  \node[ent=hoc]  (d) at (0,0)   {Drone (share $a$)};
  \node[ent=vehc] (c) at (6.6,0) {Chip in ECRV (share $g$)};
  \node[ent=uavc] (v) at (13.2,0) {UAV / neighbour};
  \foreach \n in {d,c,v}{\draw[life] (\n.south) -- ++(0,-10.4);}
  \node[act=hoc,anchor=north] at (0,-0.85) {pick $r_A$; \;$R_A=r_AP$; \;$C_A=H(R_A)$};
  \draw[msg=hoc] (0,-1.9) -- node[lbl,above]{\ding{172}\ $m,\ \ts,\ C_A$ \;(tether)} (6.6,-1.9);
  \node[act=vehc,anchor=north] at (6.6,-2.2) {tether alive? type allowed?\\ $\ts$ fresh? one session?\\ pick $r_G$; \;$R_G=r_GP$};
  \draw[msg=vehc] (6.6,-3.8) -- node[lbl,above]{\ding{173}\ $R_G$} (0,-3.8);
  \draw[msg=hoc] (0,-4.6) -- node[lbl,above]{\ding{174}\ $R_A$ \;(chip checks $H(R_A)=C_A$)} (6.6,-4.6);
  \node[act=vehc,anchor=north] at (6.6,-4.9) {$R=R_A+R_G$; \;$c=\Ht(\dots R\dots)$\\ $v_G=r_G+c\,g$; \;log};
  \draw[msg=vehc] (6.6,-6.3) -- node[lbl,above]{\ding{175}\ $v_G$} (0,-6.3);
  \node[act=hoc,anchor=north] at (0,-6.6) {$v=r_A+c\,a+v_G$};
  \draw[msg=hoc] (0,-7.6) -- node[lbl,above]{signed beacon / reply / vote $(R,v)$ \;(radio)} (13.2,-7.6);
  \node[act=uavc,anchor=north] at (13.2,-7.9) {$vP\overset{?}{=}R+c\,\pk$\\ (unchanged check)};
  \node[draw=tether,fill=tether!8,rounded corners,font=\scriptsize\sffamily\bfseries,text=tether] at (3.3,-9.6) {two short exchanges over the wire, about 200 bytes};
\end{tikzpicture}}
\caption{Joint signing over the tether. The commitment in \ding{172} forces the drone to fix its nonce before it sees the chip's.}
\label{fig:c4flow}
\end{figure}

\section{Correctness and security}
\begin{lemma}[Correctness]\label{lem:c4corr}
$aP+gP=\pk$, and the jointly produced $(R,v)$ satisfies the ordinary C1 check $vP=R+c\,\pk$; the batch equation~\eqref{eq:bv} is unchanged.
\end{lemma}
\begin{proof}
$aP+gP=X_A+X_G+dP=X+Y+h\Ppub=\pk$ (Lemma~\ref{lem:keycons}). Then
$vP=(r_A+c\,a)P+(r_G+c\,g)P=(r_A+r_G)P+c\,(a+g)P=R+c\,\pk$.
\end{proof}

\begin{theorem}[Half a key gives no signing power]\label{thm:half}
Let $G=gP=\pk-aP$. If an adversary that knows the drone share $a$ outputs a valid signature $(R,v)$ on a message the chip never co-signed, then
$(R,\ v-c\,a)$ is a valid Schnorr signature under $G$ alone. Hence forging with the drone share is as hard as forging a Schnorr signature under an unknown
key, which in the random-oracle model is as hard as the ECDLP (Lemma~\ref{lem:fork}).
\end{theorem}
\begin{proof}
$(v-c\,a)P=vP-c\,aP=R+c\,(aP+G)-c\,aP=R+c\,G$.
\end{proof}

\begin{theorem}[Share independence]\label{thm:indep}
Because $x_G$ is uniform and secret, the chip share $g=x_G+d$ is uniform from the drone's point of view, and the drone's entire memory determines $g$ only
through the public value $G=\pk-aP$. Recovering $g$ from a captured drone is therefore the ECDLP for $G$. The TA cannot sign either: it knows $d$ but
neither $x_A$ nor $x_G$.
\end{theorem}

\begin{theorem}[Capture resistance with past sessions]\label{thm:capture}
Suppose the chip runs J1--J5 with commit-then-reveal and one session at a time. Then even an adversary that controlled the drone side of earlier signing
sessions (a software takeover before the capture) and later obtains $a$ cannot sign a message the chip did not approve, unless it solves the ECDLP, in the
random-oracle model.
\end{theorem}
\begin{proof}[Proof sketch]
Plant an ECDLP challenge as $G$ and simulate the chip without $g$: because the drone must commit to $R_A$ before seeing $R_G$, the simulator can read
$R_A$ from the random-oracle query behind $C_A$, choose $v_G$ and $c$ first, set $R_G=v_GP-c\,G$ and program $\Ht$ at $R=R_A+R_G$ -- the technique used
for three-round Schnorr multi-signatures~\cite{bnforking}. A forgery is turned into a signature under $G$ (Theorem~\ref{thm:half}) and forked
(Lemma~\ref{lem:fork}) to extract $g$.
\end{proof}

\begin{remark}[Never use the naive two-round version]
If the chip answered many sessions in parallel without the commitment, a hijacked drone could choose its nonces after seeing the chip's and combine
sessions to forge -- the ROS attack~\cite{ros21}. Use commit-then-reveal with one session at a time, as above, or the standardised FROST
protocol~\cite{frost20,rfc9591} if concurrency is needed.
\end{remark}

\begin{workedbox}[title={Example 10.1 -- the textbook curve $y^2=x^3+2x+2$ over $\mathbb F_{17}$, $P=(5,1)$, order 19~\cite{paar}}]
\textbf{Setup.} TA: $s=8$, $\Ppub=8P=(13,7)$; \;$y=11$, $Y=11P=(13,10)$; \;take $h=10$, so $h\Ppub=80P=4P=(3,1)$;
\;$d=11+8\cdot10=91\equiv\mathbf{15}$ (check: $15P=Y+h\Ppub=11P+4P$ \yes).\\
Drone $x_A=6$, $X_A=6P=(16,13)$; \;chip $x_G=7$, $X_G=7P=(0,6)$; \;$X=13P=(16,4)$.\\
Shares: $\mathbf{a=6}$, \;$\mathbf{g=7+15=22\equiv3}$. \;$\pk=X+Y+h\Ppub=13P+11P+4P=28P\equiv9P=(7,6)$; nobody holds $a+g=9$.\\[4pt]
\textbf{Joint signature.} $r_A=5$: $R_A=5P=(9,16)$; \;$r_G=12$: $R_G=12P=(0,11)$; \;$R=17P=(6,14)$; \;take $c=2$.
$v_A=5+2\cdot6=17$, \;$v_G=12+2\cdot3=18$, \;$v=35\equiv\mathbf{16}$. Check: $vP=16P=(10,11)$ and $R+c\,\pk=17P+18P=35P\equiv16P$ \yes.\\[4pt]
\textbf{Thief with only $a=6$.} $r'=3$: $R'=3P=(10,6)$; take $c'=5$; $v'=3+5\cdot6=33\equiv14$, so $v'P=14P=(9,1)$. Needed: $R'+c'\pk=3P+45P\equiv10P=(7,11)$.
$14P\neq10P$ \no\ \textbf{rejected}. The missing piece is $c'gP=15P=(3,16)$, which only the chip can add.
\end{workedbox}

\section{When the tether is cut}
\begin{figure}[H]
\centering
\resizebox{\linewidth}{!}{%
\begin{tikzpicture}[font=\sffamily\small]
  \draw[-{Stealth[length=3mm]},very thick,black!60] (0,0) -- (15.4,0) node[right]{time};
  \foreach \x/\t/\c in {0.8/{$t=0$}/bad,4.6/{$\approx0$}/vehc,8.6/{$\le\delta$}/uavc,12.6/{then}/good}{
     \fill[\c] (\x,0) circle (0.14); \node[below=4pt,font=\scriptsize\sffamily\bfseries,text=\c] at (\x,0) {\t};}
  \node[card=bad,text width=3.3cm,anchor=south] at (0.8,0.45) {tether cut; drone lands or is taken\\[2pt] \textbf{chip refuses for ever}};
  \node[card=vehc,text width=3.4cm,anchor=south] at (4.6,0.45) {ECRV broadcasts the stored \textbf{kill notice}; neighbours drop the TUAV};
  \node[card=uavc,text width=3.4cm,anchor=south] at (8.6,0.45) {no new beacon can be signed; old ones expire after $\delta$; members leave};
  \node[card=good,text width=3.4cm,anchor=south] at (12.6,0.45) {members re-join a neighbour by the \textbf{full join} (fresh ECDH); old keys useless};
\end{tikzpicture}}
\caption{What happens after a cut: the thief holds half a key and an old group key that dies at the next rekey.}
\label{fig:cut}
\end{figure}

\section{Cost with the base paper's constants}
Today one TUAV signature costs $T_{SM\cdot E}+T_{SH}=0.903$\,ms. With the split, the drone does $T_{SM\cdot E}+T_{PA\cdot E}+2T_{SH}=0.915$\,ms (only
$+0.012$\,ms) and the chip $T_{SM\cdot E}+2T_{SH}=0.904$\,ms \emph{in parallel}. The tether carries about 200 bytes per signature ($C_A$, $R_A$, $R_G$,
$v_G$ and $m$), i.e.\ 1.6\,ms on a slow 1\,Mbit/s link. UAVs and neighbouring TUAVs need \textbf{no change at all}. The ECRV needs no data centre: a
32-byte share in a SIM-sized chip and one scalar multiplication per TUAV signature.

\section{What is known and what is new}
Splitting a key between a device and a helper that can switch it off is \emph{mediated cryptography}~\cite{bdtw01}; it was combined with certificateless
keys in~\cite{cbn06}, a pairing-free mediated certificateless signature appears in~\cite{yap07}, and two-party and threshold Schnorr signing is
standardised as FROST~\cite{frost20,rfc9591}. What is new is \emph{where} we place it: the helper is the ECRV at the other end of the physical tether, the
off-switch is the tether itself, and a pre-signed kill notice informs the neighbours. We found no prior work, among the papers citing the base scheme or in
the UAV authentication literature, that splits a tethered drone's key across its tether.

\begin{limitbox}
\begin{itemize}[leftmargin=*]
\item If the whole ECRV is seized together with the drone, both shares are lost (a crewed vehicle is much harder to seize than a fallen drone).
\item A software takeover while still tethered can still ask the chip to sign within its rules; C3 is the backup.
\item The group key in the drone's memory remains exposed until members move to a neighbour and rekey.
\item C4 needs a tamper-resistant chip and the careful protocol (commit-then-reveal or FROST); a full formal model of J1--J5 is future work.
\end{itemize}
\end{limitbox}
